% `tex.web` §277 `eq_save`, for the `\skip` and `\muskip` banks.
%
%     if xeq_level[p]<>cur_level then
%       begin eq_save(p,xeq_level[p]); xeq_level[p]:=cur_level;
%       end;
%
% A register is saved the *first* time it is assigned in a group and not
% again, because the group only ever has to put back the one value it found.
% The `\count` and `\dimen` banks already worked this way here; these two did
% not — every assignment pushed an entry. So `{\skip0=1pt \skip0=2pt}` wrote
% two `{restoring …}` lines where TeX writes one, and a loop that sets a
% scratch skip inside a group carried an entry per turn on the save stack.
%
% §283's `\global` is the other half, and getting it wrong is worse than a
% duplicated line: the entry stays on the stack with its level set to the
% outermost, and `unsave` *discards* it — tracing `{retaining …}`, which is how
% a reader sees that a `\global` inside a group beat the group. Deleting the
% entry instead says nothing; applying it would undo the `\global`.
\catcode`\{=1 \catcode`\}=2
\tracingonline=1 \tracingrestores=1
\skip0=1pt \muskip0=1mu \skip1=1pt
\showthe\escapechar

% Three assignments to one register, and two to another: two lines, not five.
{\skip0=2pt \skip0=3pt \skip0=4pt \muskip0=2mu \muskip0=3mu}
\showthe\skip0

% A `\global` in the middle: the group retains it rather than restoring it,
% and the value survives.
{\skip1=2pt \global\skip1=9pt \skip1=3pt}
\showthe\skip1

% Nested groups each save once, and the inner one's restore is the outer
% one's starting point.
{{\skip0=5pt}\skip0=6pt}
\showthe\skip0
\showthe\escapechar
\end
